\documentclass[11pt,a4paper]{article}

\usepackage[utf8]{inputenc}
\usepackage[T1]{fontenc}
\usepackage{lmodern}
\usepackage[margin=1in]{geometry}
\usepackage{xcolor}
\usepackage{graphicx}
\usepackage{enumitem}
\usepackage{fancyhdr}
\usepackage{booktabs}
\usepackage{titlesec}
\usepackage[colorlinks=true, linkcolor=primaryblue, urlcolor=primaryblue, citecolor=primaryblue]{hyperref}

% Define blue color scheme
\definecolor{primaryblue}{RGB}{0, 75, 150}
\definecolor{secondaryblue}{RGB}{45, 115, 190}
\definecolor{lightgray}{RGB}{240, 240, 240}

% Title Formatting
\titleformat{\section}{\Large\bfseries\color{primaryblue}}{\thesection}{1em}{}
\titleformat{\subsection}{\large\bfseries\color{secondaryblue}}{\thesubsection}{1em}{}
\titleformat{\subsubsection}{\normalsize\bfseries\color{secondaryblue}}{\thesubsubsection}{1em}{}

% Header and Footer
\pagestyle{fancy}
\fancyhf{}
\renewcommand{\headrulewidth}{0.4pt}
\renewcommand{\headrule}{\hbox to\headwidth{\color{primaryblue}\leaders\hrule height \headrulewidth\hfill}}
\fancyhead[L]{\textcolor{primaryblue}{\textbf{Adaptive Quiz Platform}}}
\fancyhead[R]{\textcolor{primaryblue}{Comprehensive Project Report}}
\fancyfoot[C]{\thepage}

\title{\vspace{-2cm}\color{primaryblue}\Huge\textbf{Adaptive Quiz Platform}\\ \Large\vspace{0.5cm}\color{secondaryblue}Comprehensive Methodological, Architectural, and Data Science Report}
\author{\textbf{Project Development Team}}
\date{\today}

\begin{document}

\maketitle
\thispagestyle{fancy}

\begin{abstract}
\noindent This document serves as the definitive, extensive logic review of the Adaptive Quiz Platform's machine learning pipeline. It is designed to provide stakeholders with a comprehensive understanding of every architectural, pedagogical, and mathematical decision made across the 9 stages of the project. Abstracting away raw code, it focuses on the strategic data-driven reasoning behind our target class choices, the curation of the question bank, the generation of synthetic student models, the prevention of data leakage, and the robust validation against four distinct real-world datasets (Data Mining, OS/ITE, Lectu, and EdNet). This report validates the system's operational readiness and highlights the sophisticated engineering required to evaluate student knowledge dynamically.
\end{abstract}

\newpage
\tableofcontents
\newpage

\section{Executive Summary and Pipeline Architecture}
The goal of the Adaptive Quiz Platform is to replace static, one-size-fits-all quizzes with an active-learning system that adapts to each student in real time. Instead of asking every student the same 20--30 questions, the system observes early performance and dynamically selects the most informative next question to rapidly classify the student's proficiency level. 

Our dual-optimisation problem consists of:
\begin{enumerate}
    \item \textbf{Level Estimation:} Given a student's partial response history, predict their true underlying knowledge level.
    \item \textbf{Question Selection:} Given the current level uncertainty, select the next question that maximises information gain.
\end{enumerate}

To solve this, we constructed a 9-stage pipeline progressing from raw question curation, through synthetic student simulation, feature engineering, predictive modelling, active learning simulation, and finally, rigorous validation against real-world educational datasets (EdNet, ENSIA Lectu). The system successfully demonstrates a \textbf{33\% reduction} in questions required to reach a stable proficiency classification when compared to non-adaptive (random) questioning.

\section{Target Class Architecture: The Rule of Three}
A foundational decision in our machine learning pipeline was defining the output space of the model. We explicitly chose to classify students into exactly three levels: \textbf{Beginner}, \textbf{Intermediate}, and \textbf{Advanced}.

\subsection{Why Not More Classes?}
It is tempting to create a highly granular scale (e.g., a 10-point scale). However, increasing the number of classes requires exponentially more data to define the narrow boundaries between them. More importantly, from a pedagogical standpoint, a 10-level system does not translate into actionable feedback during a short quiz session. If a student is a "Level 6" versus a "Level 7," the system's operational reaction remains identical. Granularity introduces misclassification noise without adding actionable value.

\subsection{Why Not Fewer Classes?}
Conversely, a binary "Pass/Fail" system is too restrictive. It fails to facilitate a truly adaptive testing environment. To measure a student's precise edge of competence, the system needs a middle ground (Intermediate) where medium-difficulty questions can be served to discern whether the student is guessing or genuinely understands the core concepts. The three-tier system flawlessly maps to standard educational item difficulties (Easy, Medium, Hard), allowing unambiguous routing decisions.

\section{Question Bank Curation and Difficulty Labeling}
The raw material of the system is the question bank. We utilized Sanfoundry, a public repository spanning multiple computer science modules. 

\subsection{The Three-Pass Filtering System}
Raw scraped data is inherently noisy. We applied three rigorous filters:
\begin{itemize}[label=\textcolor{primaryblue}{\textbullet}]
    \item \textbf{Choice Completeness:} Questions with fewer than two valid choices were discarded as scraping artifacts.
    \item \textbf{Near-Duplicate Removal:} Using TF-IDF text similarity, pairs with >80\% similarity within the same module were removed to prevent the model from memorising surface patterns instead of true difficulty signals.
    \item \textbf{Minimum Module Size:} Only modules retaining at least 150 clean questions were kept, ensuring enough depth for adaptive selection without exhausting the bank.
\end{itemize}
This yielded a pristine bank of 8,561 questions across 10 modules.

\subsection{The Difficulty Labeling Problem}
Assigning difficulty to questions is the most consequential step in the pipeline. If a "hard" question is mislabeled as "medium," the model's feature engineering collapses.

Initially, simple tertile splitting (bottom 33\% success rate = hard, top 33\% = easy) was considered. However, this suffers from \textbf{Population Confounding}: if a class happens to be extremely gifted, even hard questions will have high success rates and be mislabeled as "easy." 

To solve this, we implemented \textbf{Proficiency-Weighted Wrong Rate (PWWR)}. This means a wrong answer from an historically strong student carries much more weight in classifying a question as "hard" than a wrong answer from a struggling student. We mapped PWWR to absolute fixed thresholds (e.g., PWWR $\leq$ 0.30 = easy), completely decoupling question difficulty from the varying abilities of the specific students who take them.

\section{The Synthetic Population Model}
Machine learning requires data, but as a new platform, we faced a "cold start" problem with zero existing adaptive user logs. To build and validate the system, we required a synthetic population heavily grounded in Educational Data Mining (EDM) theory.

\subsection{Population Distribution: N(0.55, 0.18)}
We generated 1,000 synthetic students with latent proficiency scores drawn from a Gaussian distribution with mean 0.55 and standard deviation 0.18. Why this specific distribution?
\begin{itemize}[label=\textcolor{primaryblue}{\textbullet}]
    \item A uniform distribution artificially inflates model accuracy because there are too many students at the extreme ends (0.1 or 0.9 proficiency) who are trivially easy to classify.
    \item The N(0.55, 0.18) curve places the level boundaries (0.45 and 0.70) exactly inside the highest-density regions of the curve. This forces the model to learn from thousands of highly ambiguous, borderline students—the exact scenarios where adaptive systems provide the most value.
\end{itemize}

\subsection{Modeling Guessing and Slipping}
Borrowing from Item Response Theory (IRT), we did not equate true proficiency directly to quiz score. We introduced a \textbf{Guessing Parameter} (weak students randomly guessing correctly) and a \textbf{Slipping Parameter} (advanced students making careless mistakes). Response times were generated via an inverse proficiency model (struggling students take longer).

We deliberately injected a dedicated cohort of 151 "Borderline Students" who sit exactly ±0.02 from the class boundaries to serve as the ultimate stress test for our adaptive algorithms.

\section{Feature Engineering: The 25-Feature Vector}
A raw quiz score (e.g., 7/10) is a poor indicator of true mastery because it cannot distinguish a lucky guesser from a fluent learner. We engineered a 25-dimensional behavioral vector extracted at multiple checkpoints (k=5, 10, 15... questions).

\subsection{Categorical Breakdown}
\begin{itemize}[label=\textcolor{primaryblue}{\textbullet}]
    \item \textbf{Difficulty-Stratified Accuracy (7 features):} \texttt{acc\_easy}, \texttt{acc\_medium}, \texttt{acc\_hard}, and boolean flags to indicate if a difficulty tier has been seen at all. A student who aces easy but fails hard questions is fundamentally different from a student who fails both.
    \item \textbf{Temporal Momentum (4 features):} \texttt{last3\_acc}, \texttt{error\_streak}, \texttt{correct\_streak}. These track whether a student is improving or plateauing mid-session.
    \item \textbf{Response Time Dynamics (2 features):} \texttt{avg\_time}, \texttt{time\_trend}. Fast answers on complex problems often signal guessing.
    \item \textbf{Per-Module Knowledge (12 features):} Accuracy and visibility flags for 6 distinct computer science modules, tracking cross-domain breadth.
\end{itemize}

\subsection{The Time-Feature Artifact Warning}
Our ablation studies revealed that \texttt{avg\_time} accounted for 64\% of total permutation importance. However, we explicitly flagged this as a \textit{synthetic artifact}. In our synthetic generator, time was derived directly from proficiency. In the real world (EdNet, ENSIA), time is far noisier due to reading speed and platform latency. We verified via ablation that removing time entirely only slightly impacts the core accuracy, proving the 25-feature vector is structurally sound independently of response times.

\section{Classifier Selection and Calibration}
We evaluated Logistic Regression, Gaussian Naive Bayes, Decision Trees, and Random Forests on the 25-feature matrix. 

\subsection{Why Logistic Regression Triumphed over Random Forest}
Random Forest achieved 88.77\% accuracy, while Logistic Regression achieved 89.20\%. But accuracy was not the primary deciding factor. 

The adaptive algorithms (Uncertainty Sampling, Entropy) require \textbf{Calibrated Probabilities}. If a model says it is "90\% confident" a student is Advanced, it must statistically be correct 9 times out of 10 in reality. Random Forest trees vote in fractions that are notoriously uncalibrated. Logistic Regression, utilizing a sigmoid/softmax function, produces highly calibrated probabilities naturally. 

To address the population imbalance (more Intermediate than Beginner/Advanced), we utilized \texttt{class\_weight='balanced'}, trading 1-2\% overall accuracy for a massive +14\% jump in Beginner recall—ensuring struggling students are never left behind.

\section{Data Leakage Prevention: Grouped Cross-Validation}
Machine learning models in education are incredibly susceptible to "Data Leakage"—where a model memorizes a student rather than learning a general rule.

\subsection{The Perils of Standard K-Fold}
Because we extract feature vectors at checkpoints (k=5, k=10, k=15) for a single student session, standard Stratified K-Fold CV randomly splits these rows. A model might train on a student's k=5 row, and test on the exact same student's k=10 row. 
When we tested standard CV, our accuracy falsely inflated to 92\%. 

\subsection{The Grouped CV Fix}
We implemented \textbf{Grouped Cross-Validation (groups=student\_id)}. This forces all rows belonging to a single student into the exact same fold. The model is forced to evaluate students it has genuinely never seen before, dropping accuracy back to an honest, rigorous 89.20\%.

Furthermore, any dataset-wide metrics (like PWWR difficulty labels or standard scalers) were strictly calculated \textit{inside} the training folds to ensure zero future information leaked into the test fold.

\section{Adaptive Question Selection Strategies}
With the classifier established, we simulated four adaptive selection strategies against a random non-adaptive baseline.

\begin{enumerate}
    \item \textbf{Random Baseline:} Selects uniformly.
    \item \textbf{Uncertainty Sampling:} Selects the question that forces the model to decide between the two most likely classes.
    \item \textbf{Entropy Reduction:} Selects the question that mathematically minimises expected future entropy.
    \item \textbf{Query By Committee (QBC):} An ensemble of 4 models vote. The system selects the question that causes the most disagreement among the models.
\end{enumerate}

\subsection{The 33\% Efficiency Gain}
The QBC strategy emerged as the definitive winner. By seeking maximum committee disagreement, it reached the 80\% population accuracy target by question 6. The random baseline required 9 questions to achieve the same confidence. Over a 30-question session, this represents a **33\% reduction** in assessment overhead, directly returning time to actual student learning.

\section{Real Data Validation: Cross-Domain Generalization}
A synthetic model is only theoretical until validated on real humans. We integrated four real-world datasets to test generalizability.

\subsection{EdNet and Lectu (ENSIA)}
We tested the pipeline on 1,107 students from EdNet KT1 (a massive Korean educational dataset) and 118 students from the Lectu platform at ENSIA. 

The grouped, leakage-free Cross-Validation results were staggering:
\begin{itemize}[label=\textcolor{primaryblue}{\textbullet}]
    \item \textbf{Lectu/ENSIA:} 96.6\% Accuracy
    \item \textbf{EdNet KT1:} 96.9\% Accuracy
    \item \textbf{EdNet KT2:} 96.7\% Accuracy
\end{itemize}

\subsection{The Operational Definition of Level}
Critics might ask: Is 96.9\% accuracy "too high"? Does it indicate circularity? 
The answer is no. In our platform, "Level" is an \textit{operational definition} derived directly from quiz performance (e.g., >70\% hard accuracy = Advanced). The ML model is not predicting a mystical hidden state; it is learning our operational mapping and applying it seamlessly to behavioral feature vectors of previously unseen students. The high CV proves that human student behavior cleanly stratifies along these operational boundaries.

\section{Real Data Anomalies: The Online Inflation Problem}
Not all real data is clean. We encountered critical pedagogical anomalies during validation.

\subsection{Data Mining (DM) Quizzes}
For a Data Mining module, we compared students' unsupervised online Google Form quiz scores against their highly supervised, in-person written exams.
\begin{itemize}[label=\textcolor{primaryblue}{\textbullet}]
    \item \textbf{In-person Exam Mean:} 56.1\%
    \item \textbf{Online Quiz Mean:} 83.2\%
\end{itemize}
This +27.2\% gap highlights massive online score inflation. The correlation between quiz and exam was only $r = 0.14$. If we try to train the ML model to predict Exam Grades using Online Quiz features, accuracy drops to ~40\%. This proves a fundamental theorem for the platform: \textbf{Feature-Label Consistency}. The pipeline only functions accurately when the environment generating the behavioral features (the online app) is under the same supervision conditions as the environment assigning the true labels.

\subsection{OS and ITE Quizzes}
A dataset of OS and ITE quizzes showed a 92\% class imbalance, where 578 out of 626 students were classified as "Advanced." This was caused by severe self-selection bias (only excellent students submitted their logs) and unsupervised at-home inflation. A model trained here would trivially predict "Advanced" for everyone. This reinforces the necessity of our balanced synthetic population generation to train the core platform engine.

\section{Conclusion}
The Adaptive Quiz Platform represents a sophisticated intersection of Educational Data Mining and Machine Learning. By rigorously curating data, leveraging Item Response Theory for synthetic populations, engineerings behavioral features immune to raw-score luck, and actively preventing data leakage, the system achieves highly calibrated, predictive accuracy. 

The transition from static to adaptive questioning cuts assessment time by 33\%, and its successful validation on datasets like EdNet proves the platform is robust, scientifically defensible, and ready for deployment.

\end{document}
